\section{Introduction}
\label{sec:intro}

Differentiable rendering has made it possible to optimize an explicit scene representation directly from posed images. Gaussian splatting~\cite{kerbl2023gaussians} is the most used such representation today, but a growing family of methods instead optimizes an \emph{opaque triangle mesh}~\cite{meshsplatting2025,milo2025,rade2024,gof2024,bog2024}: a set of triangles with per-vertex position, color and view-dependent appearance, rendered by an ordinary rasterizer. The appeal is practical. A mesh renders in any engine, on a phone, with depth, collision and relighting available for free. The cost is size: a trained \MS scene occupies 430--780~MB in 32-bit floats, two to three orders of magnitude more than a compressed Gaussian scene~\cite{hacpp2025,compgs2024}, and no compression scheme exists for this representation.

The obvious route, train first and simplify afterwards, does not work for opaque meshes. Every pixel is covered by exactly one triangle, so the vertices of a trained mesh are placed and colored for the neighbours they have; when a decimator removes some of them, the survivors are left where they are, and fine-tuning does not move them back. We measure this directly (Sec.~\ref{sec:e1}): at the same vertex count, a mesh trained to that size from the start is better than a trained-then-simplified mesh on 62 of 63 scene--budget pairs, by 1.85\dB on average and by up to 5.5\dB at the smallest budgets. With Gaussians the opposite holds; pruning a trained splat scene by half costs a fraction of a decibel~\cite{lightgaussian2024,compact3dgs2024}. This is the central observation of the paper: for opaque meshes, the size of the file has to be decided \emph{while the mesh grows}, when the optimizer can still adapt every vertex to the triangles it will end up with.

\BM turns this observation into a system. It keeps \MS's optimizer and schedule and changes only how the mesh grows and shrinks. (i)~\textbf{Rate-aware growth.} \MS densifies by splitting triangles sampled by an importance score. \BM instead accepts a split only when the rendering error it removes is worth the bytes it adds to the \emph{encoded} file, with the bytes measured by the actual encoder on the current mesh. (ii)~\textbf{Size as an input.} The user gives a file size in megabytes. A controller plans how many vertices that size buys, sets the price of a byte at every growth step so that the plan is followed, and in the second half of training applies the same test in reverse, collapsing the edges whose error per byte saved is lowest, until the measured file size meets the target. (iii)~\textbf{A codec with a standalone decoder.} Positions are stored on a per-scene grid, connectivity with a context-based arithmetic coder, and the 45 view-dependent coefficients of each vertex as two indices into per-scene dictionaries. A decoder with no access to the training data reconstructs a standard indexed mesh in a few seconds on one CPU core.

We evaluate on Mip-NeRF~360, Tanks and Temples, NeRF-Synthetic and Deep Blending, the last one untouched during development. At the same vertex count, and before any entropy coding, rate-aware growth is better than \MS trained with a vertex cap on 44 of 45 pairs (+0.27\dB on real scenes). Given a target size, 83 of 89 files land within $\pm10\%$ of it with one setting per capture class, at a cost of 0.12\dB relative to the best the same growth rule reaches at that size when the size is not prescribed. At 20, 10 and 5~MB the files are 31$\times$, 60$\times$ and 122$\times$ smaller than the FP32 model at $-0.30$, $-0.66$ and $-1.32$\dB; against the same meshes through Draco or MPEG with 8-bit appearance, our files are 3--18$\times$ smaller. We also report what the method does not do: at the same file size an opaque mesh remains 3--4\dB below the best Gaussian codecs, a gap of the representation rather than of the coder, and we state it with the numbers.

\noindent\textbf{Contributions.}
\begin{itemize}\setlength\itemsep{0pt}
\item A rate-aware growth rule for differentiable opaque meshes that prices each split in encoded bytes, together with the finding that such meshes must be sized during training: trained-to-size beats trained-then-simplified on 62/63 pairs by 1.85\dB.
\item A size controller that takes the file size as input, plans vertex growth from byte costs measured by the encoder during training, and trims by the same rule, landing within $\pm10\%$ on 83/89 runs across four datasets with one setting per capture class.
\item A codec and standalone decoder for this representation, 32--125$\times$ smaller than FP32 at 0.3--1.3\dB, with an analysis of where the bytes go: connectivity is 40\% of the file and lossless; the shared appearance dictionary carries almost all of the loss.
\end{itemize}
